% !TEX root = scientific-computing.tex

\chapter{Managing Packages in Julia}

Although we can use the \texttt{Pkg} package to handle packages, this
section will use the terminal REPL to handle any non-standard packages.
Additional documentation on this is given in
\href{Julia's\%20package\%20manager\%20help\%20pages}{https://docs.julialang.org/en/latest/stdlib/Pkg/}
First, open up a terminal version of julia (generally by opening the
application that you downloaded). You will get:

\begin{minted}{julia}
              _
   _      _ _(_)_      |  Documentation: https://docs.julialang.org
  (_)    | (_) (_)     |
  _ _   _| |_  __ _    |  Type "?" for help, "]?" for Pkg help.
  | | | | | | |/ _` |  |
  | | |_| | | | (_| |  |  Version 1.0.0 (2018-08-08)
 _/ |\__'_|_|_|\__'_|  |  Official https://julialang.org/ release
|__/                   |
\end{minted}

or similar and then

\begin{minted}{julia}
julia>
\end{minted}

which means we're ready to handle julia commands. If we type
\texttt{{]}}, then the prompt turns into:

\begin{minted}{julia}
(v1.0) pkg>
\end{minted}

There are a number of commands that we will cover here:

\begin{itemize}

\item
  add
\item
  remove (rm)
\item
  develop (dev)
\item
  status
\item
  update (up)
\item
  test
\item
  build
\item
  precompile
\end{itemize}

\subsection{Adding a package}

In the package command line, type \texttt{add} \emph{package\_name} to
add the package. For example, to add the \texttt{ForwardDff} package:

\begin{minted}{julia}
add ForwardDiff
\end{minted}

and it is not installed, you will get something like:

\begin{verbatim}
Updating registry at `~/.julia/registries/General`
  Updating git-repo `https://github.com/JuliaRegistries/General.git`
  Updating `~/.julia/environments/v1.0/Project.toml`
  [f6369f11] + ForwardDiff v0.9.0
  Updating `~/.julia/environments/v1.0/Manifest.toml`
  [9e28174c] + BinDeps v0.8.10
  [bbf7d656] + CommonSubexpressions v0.2.0
  [163ba53b] + DiffResults v0.0.3
  [b552c78f] + DiffRules v0.0.7
  [f6369f11] + ForwardDiff v0.9.0
  [77ba4419] + NaNMath v0.3.2
  [276daf66] + SpecialFunctions v0.7.0
  [90137ffa] + StaticArrays v0.8.3
\end{verbatim}

A few things to note:

\begin{itemize}

\item
  The line after the 2nd \texttt{Updating} line is the package (and
  version) that you are installing.
\item
  All of the lines after the 3rd \texttt{Updating} line is all of the
  packages that this depends on.
\item
  The + sign means that the package is being added.
\end{itemize}

If you want to add multiple packages at the same time, say packages A, B
and C, type \texttt{add\ A\ B\ C}. You can also add particular versions
of a package (often for testing or to avoid a bug). For example, if you
want version 0.3.0 of \texttt{ForwardDiff} type:

\begin{verbatim}
add ForwardDiff@0.3.0
\end{verbatim}

You will then get info on the dependencies on that version.

\hypertarget{package-status}{%
\subsubsection{Package Status}\label{package-status}}

The \texttt{status} command (or \texttt{st}) will just list all of the
main packages installed. For example,

\begin{minted}{julia}
Status `~/.julia/environments/v1.0/Project.toml`
[c52e3926] Atom v0.7.6
[336ed68f] CSV v0.3.1
[a93c6f00] DataFrames v0.13.1
[f6369f11] ForwardDiff v0.9.0
[7073ff75] IJulia v1.12.0
[e5e0dc1b] Juno v0.5.3
[1a8c2f83] Query v0.10.0
\end{minted}

and note that these are just the packages added by the \texttt{add}
command, not all of the dependencies. If you want all of the
dependencies as well, type \texttt{st\ -\/-manifest} and I get a huge
list of packages.

\hypertarget{removing-a-package}{%
\subsubsection{Removing a Package}\label{removing-a-package}}

You can remove a package by typing \texttt{remove} or \texttt{rm} then
the package name. If I want to remove the \texttt{ForwardDiff} package,
then

\begin{minted}{julia}
remove ForwardDiff
\end{minted}

we get the following:

\begin{minted}{julia}
Updating `~/.julia/environments/v1.0/Project.toml`
 [f6369f11] - ForwardDiff v0.3.0
 Updating `~/.julia/environments/v1.0/Manifest.toml`
 [49dc2e85] - Calculus v0.4.1
 [c5cfe0b6] - DiffBase v0.2.0
 [f6369f11] - ForwardDiff v0.3.0
 [77ba4419] - NaNMath v0.3.2
\end{minted}

Note:

\begin{itemize}

\item
  The \texttt{rm} command removes the package from the list of available
  packages, but doesn't remove them from your harddrive.
\item
  If you want to see everything installed, navigate to the
  \texttt{\textasciitilde{}/.julia/packages} directory, which is where
  they are stored.
\end{itemize}

\hypertarget{updating-packages}{%
\subsubsection{Updating packages}\label{updating-packages}}

If you type \texttt{update} or \texttt{up} you will update all of the
installed packages (and dependencies). For example:

\begin{minted}{julia}
Updating `~/.julia/environments/v1.0/Project.toml`
 [7073ff75] ↑ IJulia v1.11.1 ⇒ v1.12.0
 Updating `~/.julia/environments/v1.0/Manifest.toml`
 [7073ff75] ↑ IJulia v1.11.1 ⇒ v1.12.0
 [b85f4697] ↑ SoftGlobalScope v1.0.5 ⇒ v1.0.7
 [5e66a065] ↑ TableShowUtils v0.1.1 ⇒ v0.2.0
\end{minted}

and all updates will be with an ↑. If you only want to update a single
package, type the name after \texttt{update}.

\hypertarget{building-packages}{%
\subsubsection{Building Packages}\label{building-packages}}

Generally a package is built after it is installed. Building a package
might include running code (or unpacking files) after it is installed.
Sometimes if things get wonky, rebuilding is a good thing to do.

\begin{minted}{julia}
build
\end{minted}

or if you only want say \texttt{IJulia} built,

\begin{minted}{julia}
build IJulia
\end{minted}

\subsubsection{Precompiling packages}

When a package is used, often it requests to be compiled. For example,
when \texttt{using\ Primes}, then following is shown:

\begin{minted}{julia}
[ Info: Precompiling Primes [27ebfcd6-29c5-5fa9-bf4b-fb8fc14df3ae]
\end{minted}

and basically some code is compiled beforehand, generally to speed up
code. You can precompile all code with

\begin{minted}{julia}
precompile
\end{minted}

and it may take a while, but you won't have to wait, when you load the
package with the \texttt{using} command.

\hypertarget{testing-a-packages}{%
\subsubsection{Testing a Packages}\label{testing-a-packages}}

To test a packge, say the \texttt{ForwardDiff} package, then

\begin{minted}{julia}
test ForwardDiff
\end{minted}

It list all of the dependencies first, and then runs a number of tests
(and we will show how to write tests soon) and timing information. After
a while, it finishes sucessfully.
